% ---------------------------------------------------------------------------
% bibliography.tex -- a reading list rather than a citation apparatus.
% ---------------------------------------------------------------------------
\cleardoublepage
\phantomsection
\addcontentsline{toc}{chapter}{Further Reading}
\chapter*{Further Reading}
\markboth{Further Reading}{Further Reading}

\section*{Standards}

\begin{itemize}
  \item IEEE Std 1800-2017, \emph{IEEE Standard for SystemVerilog --- Unified
        Hardware Design, Specification, and Verification Language}. The
        language reference manual. Chapter~9 (processes), Chapter~10
        (assignment statements) and Chapter~4 (scheduling semantics) are the
        three a VHDL designer should read first.
  \item IEEE Std 1076-2008, \emph{IEEE Standard VHDL Language Reference
        Manual}, and IEEE Std 1076-2019 for the newer features discussed in
        \refdesignbook.
  \item IEEE Std 1800.2-2020, \emph{Universal Verification Methodology
        Language Reference Manual}.
  \item IEEE Std 1364-2005, the last standalone Verilog standard, useful only
        when reading legacy code.
\end{itemize}

\section*{Books}

\begin{itemize}
  \item C.~Spear and G.~Tumbush, \emph{SystemVerilog for Verification},
        3rd~ed., Springer. The standard text on the class-based testbench,
        and the natural sequel to \cref{ch:tb}.
  \item S.~Sutherland, S.~Davidmann and P.~Flake, \emph{SystemVerilog for
        Design}, 2nd~ed., Springer. The design-side counterpart, written by
        people who were in the room when the language was standardised.
  \item D.~Thomas, \emph{The Verilog Hardware Description Language}, and
        P.~Ashenden, \emph{The Designer's Guide to VHDL}. The two classics,
        for when you want to know why something is the way it is.
  \item P.~Ashenden and J.~Lewis, \emph{VHDL-2008: Just the New Stuff}.
        Short, and it covers most of what \refdesignbook says VHDL can do.
  \item M.~Keating, \emph{The Simple Art of SoC Design}. Not about either
        language, and worth reading anyway.
\end{itemize}

\section*{Papers and application notes}

\begin{itemize}
  \item C.~Cummings, \emph{Nonblocking Assignments in Verilog Synthesis, Coding
        Styles That Kill!}, SNUG 2000. The clearest published explanation of
        the blocking/non-blocking rules and the races they exist to prevent.
        Everything else by the same author on \code{always\_comb},
        \code{unique}/\code{priority} and clock-domain crossing is worth the
        time.
  \item J.~E.~Volder, \emph{The CORDIC Trigonometric Computing Technique},
        IRE Transactions on Electronic Computers, 1959. The original paper
        behind \refexamplebook.
  \item J.~S.~Walther, \emph{A Unified Algorithm for Elementary Functions},
        Spring Joint Computer Conference, 1971. The generalisation.
  \item R.~Andraka, \emph{A Survey of CORDIC Algorithms for FPGA Based
        Computers}, FPGA~1998. The practical implementation reference,
        including the gain, the convergence range and the quadrant folding
        used in \refexamplebook.
\end{itemize}

\section*{Tools and online resources}

\begin{itemize}
  \item \emph{Verilator} manual, \url{https://verilator.org/guide/}. Read the
        chapters on the language subset and on the C++ API before writing a
        C++ testbench.
  \item \emph{cocotb} documentation, \url{https://docs.cocotb.org/}.
  \item \emph{GTKWave} (\url{https://gtkwave.sourceforge.net}) and
        \emph{Surfer} (\url{https://surfer-project.org}), the two waveform
        viewers that read VCD and FST.
  \item \emph{OSVVM} (\url{https://osvvm.org}) and \emph{UVVM}
        (\url{https://uvvm.org}), the two VHDL verification frameworks
        referred to throughout \refdesignbook and \cref{ch:uvm}.
  \item The \emph{lowRISC} and \emph{OpenTitan} SystemVerilog style guide,
        \url{https://github.com/lowRISC/style-guides}. The best free
        statement of what disciplined SystemVerilog looks like, and the
        source of the naming conventions used in this book.
  \item \emph{Verible}, \url{https://github.com/chipsalliance/verible}, for
        formatting and linting.
\end{itemize}
